% Template Báo cáo BTL số 2 theo chuẩn nhà xuất bản Elsevier (cas-sc)
\documentclass[a4paper,fleqn]{cas-sc}
\usepackage[numbers]{natbib}
\usepackage{amsmath,amssymb,amsfonts}
\usepackage{graphicx}
\usepackage{booktabs}
\usepackage{hyperref}

\begin{document}
\let\WriteBookmarks\relax
\def\floatpagepagefraction{1}
\def\textpagefraction{.001}

\shorttitle{SAT Encodings for the Bandwidth Coloring Problem on GEOM Instances}
\shortauthors{Group Project - Modern Issues in Computer Science}

\title [mode = title]{An Experimental Study of Pairwise and Partial-Ordering SAT Encodings for the Bandwidth Coloring Problem}

\author[1]{Thành viên 1}[role=Co-author: Pairwise SAT & AMO Encodings]
\author[1]{Thành viên 2}[role=Co-author: Order / Partial-Ordering SAT & Incremental Solving]
\author[1]{Thành viên 3}[role=Co-author: Exact ILP/CP Baselines & Visualization]
\affiliation[1]{organization={University of Engineering and Technology, Vietnam National University}, city={Hanoi}, country={Vietnam}}

\begin{abstract}
In this report, we study and evaluate Boolean Satisfiability (SAT) encodings for the Bandwidth Coloring Problem (BCP) on the benchmark suite of 33 geometric graphs (\texttt{GEOM}). Specifically, we implement and compare the assignment-based / pairwise encoding (\texttt{ASS-S-B}) equipped with various At-Most-One (AMO) cardinality constraints (Pairwise, Sequential Counter, Binary, Commander, and Product via PySAT) against the partial-ordering / order encodings (\texttt{POP-S-B} and \texttt{POPH-S-B}) proposed by Faber, Jabrayilov, and Mutzel (SAT 2024). Both multi-shot SAT and incremental SAT solving strategies are investigated and benchmarked against exact Integer Linear Programming (Gurobi, CPLEX MIP) and Constraint Programming (CPLEX CP, OR-Tools CP-SAT) formulations.
\end{abstract}

\begin{keywords}
Bandwidth Coloring Problem \sep SAT Encoding \sep Order Encoding \sep Pairwise Encoding \sep Incremental SAT \sep GEOM Benchmark
\end{keywords}

\maketitle

\section{Introduction}
% [Thành viên 1 phụ trách]

\section{Preliminaries and Problem Definition}
% [Thành viên 1 phụ trách: Định nghĩa BCP, cận trên tham lam H(G)]

\section{SAT Encodings and Solving Strategies for BCP}
\subsection{Assignment-based Encoding (\texttt{ASS-S-B}) and AMO Variants}
% [Thành viên 1 phụ trách: Công thức (8a), (9a)-(9d), (11) và 5 biến thể AMO trong PySAT]

\subsection{Partial-Ordering / Order Encodings (\texttt{POP-S-B} and \texttt{POPH-S-B})}
% [Thành viên 2 phụ trách: Công thức (12a), (12b), (16a)-(16e), (19), (20)]

\subsection{Multi-shot SAT vs. Incremental SAT Solving}
% [Thành viên 2 phụ trách: Descending linear search với Multi-shot và Incremental SAT]

\section{Exact Mathematical and Constraint Programming Baselines}
% [Thành viên 3 phụ trách: Mô hình ILP (ASS-I-B, POP-I-B, POPH-I-B) và CP / CP-SAT]

\section{Experimental Evaluation on 33 GEOM Instances}
% [Thành viên 3 tổng hợp bảng Table 3, Table 5 và biểu đồ từ thư viện cnsplots]

\section{Conclusion}
% [Cả 3 thành viên]

\section*{Data Availability}
All source code, benchmark instances (\texttt{GEOM}), experimental logs, and plotting scripts are publicly available at: \url{https://github.com/dxbang1383/BCP}.

\bibliographystyle{cas-model2-names}
\bibliography{references}

\end{document}
